\documentclass[12pt, a4paper, notitlepage]{report}

\usepackage{abstract}
\usepackage[english,polish]{babel}
\usepackage[utf8]{inputenc}
\usepackage[titletoc,page]{appendix}

\usepackage[acronym, toc]{glossaries}
\makeglossaries
\newglossaryentry{agent} {name=agent, description={..}}
\newglossaryentry{complex_behaviour} {name=zachowanie złożone, description={...}}
\newglossaryentry{state_machine} {name=automat skończony, description={...}}
\newglossaryentry{effector} {name=efektor, description={...}}
\newglossaryentry{sensor} {name=sensor, description={...}}
\newglossaryentry{singleton} {name=singleton, description={...}}
\newglossaryentry{git} {name=Git, description={...}}

\newglossaryentry{udp} {name=UDP, description={User Datagram Protocol}}
\newglossaryentry{arm} {name=ARM, description={Advanced RISC Machine}}

% \newacronym{udp}{UDP}{User Datagram Protocol}
% \newacronym{arm}{ARM}{Advanced RISC Machine}

\newglossaryentry{latex}
{
    name=latex,
    description={Is a mark up language specially suited for scientific documents}
}

\newglossaryentry{maths}
{
    name=mathematics,
    description={Mathematics is what mathematicians do}
}

\newglossaryentry{formula}
{
    name=formula,
    description={A mathematical expression}
}



\newacronym{lcm}{LCM}{Least Common Multiple}


\usepackage[T1]{fontenc}
\usepackage[backend=biber,style=numeric-comp,sorting=nty,defernumbers=true]{biblatex}
\usepackage[pdftex]{graphicx}

\usepackage{hyperref}
\hypersetup{
    hidelinks,
}

\usepackage{wrapfig}
\usepackage[font=small,labelfont=bf,justification=centering]{caption}
\usepackage{subcaption}
\usepackage{listings}
 \lstset{
  basicstyle=\footnotesize\ttfamily,
  captionpos=b,
  breakatwhitespace=true,
  breaklines=true,
  frame=L,
  xleftmargin=\parindent,
 }
\usepackage{changepage}
\usepackage{pgfplots}
\pgfplotsset{compat=1.8}

\addtolength{\textwidth}{1cm}
\addtolength{\textheight}{1cm}

\overfullrule=5mm

\addbibresource{bibliography.bib}
\frenchspacing
\linespread{1.0}

\title{Paweł Kaczyński}
\author{Możliwości programowania klocka Lego EV3.}
\date{1 marca 1993 r.}

\begin{document}

\begin{titlepage}
    \begin{center}

        % Rok akademicki
        \begin{flushright}
            \begin{minipage}{0.35\textwidth}
                Rok akademicki 2015/2016
            \end{minipage}
        \end{flushright}

        % PW WEiTi
        \begin{flushleft}
            \begin{minipage}{0.57\textwidth}
                \begin{center}
                    \small
                    Politechnika Warszawska \\
                    Wydział Elektroniki i Technik Informacyjnych \\
                    Instytut Informatyki
                \end{center}
            \end{minipage}
        \end{flushleft}

        ~\\[2.0cm]

        % Tytuł i autor.
        \includegraphics[width=0.17\textwidth]{tytulowa_res/logo_pw.png}~\\[0.3cm]

        \Large PRACA DYPLOMOWA INŻYNIERSKA\\[0.5cm]

        \large Paweł Kaczyński \\[0.5cm]

        \LARGE \textbf{Możliwości programowania\\klocka LEGO EV3.}

        \vfill

        % Opiekun.
        \begin{flushright}
            \begin{minipage}{0.5\textwidth}
                \begin{center} \normalsize
                    Opiekun pracy\\
                    ~dr~inż. Henryk Dobrowolski
                \end{center}
            \end{minipage}
        \end{flushright}

        ~\\[0.1cm]

        % Komisja
        \begin{flushleft}
            \begin{minipage}{0.47\textwidth}
                \begin{flushleft} \normalsize
                    Ocena: \dotfill \\[0.7cm]
                    .\dotfill \\
                    \begin{center}
                        Podpis Przewodniczącego\\
                        Komisji Egzaminu Dyplomowego
                    \end{center}
                \end{flushleft}
            \end{minipage}
        \end{flushleft}

    \end{center}
\end{titlepage}

\begin{wrapfigure}{l}{0.30\textwidth}
  \vspace{-20pt}
  \begin{center}
    \includegraphics[width=0.30\textwidth]{tytulowa_res/logo_pw.png}
  \end{center}
  \vspace{-20pt}
\end{wrapfigure}

\noindent
\textbf{Specjalność: }Inżynieria Systemów Informacyjnych\\
\textbf{Data urodzenia: }1 marca, 1993 r.\\
\textbf{Data rozpoczęcia studiów: }1 listopada, 2012 r. \ \\\\\\\\\\\\\\

\begin{center}
    \large{\textbf{Życiorys}}
\end{center}
Lorem ipsum dolor sit amet, consectetur adipiscing elit. Proin pellentesque nisl vitae tellus tempor aliquam. Nullam consequat laoreet pretium. Vivamus vehicula, lectus eu elementum congue, magna dolor pellentesque quam, quis suscipit urna sem sit amet eros. Vivamus tincidunt, leo id egestas condimentum, diam ligula consequat libero, non malesuada orci enim non dolor. Duis nec dolor sit amet ipsum feugiat malesuada ut quis nibh. Quisque ac feugiat enim, rhoncus aliquam justo. Nulla facilisis convallis mauris id mattis. Nam vestibulum, quam non vehicula scelerisque, arcu arcu porta nulla, eget suscipit nisi mauris vel nibh. Donec condimentum, tellus bibendum semper pulvinar, dui lacus dignissim libero, sit amet aliquam justo lacus et magna.\\\\

\begin{flushright}
    \begin{minipage}{.5\linewidth}
        \dotfill
        \begin{flushright}
            Podpis studenta:
        \end{flushright}
    \end{minipage}
\end{flushright}

\vfill

\noindent
\large{\textbf{Egzamin Dyplomowy}}\\
Złożył egzamin dyplomowy w dniu ........ 2016 roku\\
z wynikiem \dotfill\\
Ogólny wynik studiów \dotfill\\
Dodatkowe wnioski i uwagi Komisji \dotfill\\
.\dotfill\\
.\dotfill

% Streszczenie w języku Polskim
% The \setlength commands I found in the abstract package are:
%
% \absleftindent
% \absrightindent
% \absparindent
% \absparseptext
\setlength{\absleftindent}{0cm}
\setlength{\absrightindent}{0cm}
\renewcommand{\abstractnamefont}{\normalfont\Large\bfseries}
\begin{abstract}
    \large
    Lorem ipsum dolor sit amet, consectetur adipiscing elit. Proin pellentesque nisl vitae tellus tempor aliquam. Nullam consequat laoreet pretium. Vivamus vehicula, lectus eu elementum congue, magna dolor pellentesque quam, quis suscipit urna sem sit amet eros. Vivamus tincidunt, leo id egestas condimentum, diam ligula consequat libero, non malesuada orci enim non dolor. Duis nec dolor sit amet ipsum feugiat malesuada ut quis nibh. Quisque ac feugiat enim, rhoncus aliquam justo. Nulla facilisis convallis mauris id mattis. Nam vestibulum, quam non vehicula scelerisque, arcu arcu porta nulla, eget suscipit nisi mauris vel nibh. Donec condimentum, tellus bibendum semper pulvinar, dui lacus dignissim libero, sit amet aliquam justo lacus et magna.
    \\

    \noindent
    Słowa kluczowe: LEGO Mindstorms EV3, agent mobilny, zachowanie.
\end{abstract}

~\\[0.1cm]
% Linia pozioma
\noindent \rule{\linewidth}{0.2mm}

% Streszczenie w języku angielskim
\begin{otherlanguage}{english}
\renewcommand{\abstractname}{}    % clear the title
\renewcommand{\absnamepos}{empty} % originally center

\begin{abstract}
    ~\\[0.05cm]
    \begin{center}
        {\Large Implementation of multiplayer, turn-based strategic game, available through web browsers}
    \end{center}
    ~\\[0.1cm]
    \large
    Lorem ipsum dolor sit amet, consectetur adipiscing elit. Proin pellentesque nisl vitae tellus tempor aliquam. Nullam consequat laoreet pretium. Vivamus vehicula, lectus eu elementum congue, magna dolor pellentesque quam, quis suscipit urna sem sit amet eros. Vivamus tincidunt, leo id egestas condimentum, diam ligula consequat libero, non malesuada orci enim non dolor. Duis nec dolor sit amet ipsum feugiat malesuada ut quis nibh. Quisque ac feugiat enim, rhoncus aliquam justo. Nulla facilisis convallis mauris id mattis. Nam vestibulum, quam non vehicula scelerisque, arcu arcu porta nulla, eget suscipit nisi mauris vel nibh. Donec condimentum, tellus bibendum semper pulvinar, dui lacus dignissim libero, sit amet aliquam justo lacus et magna.\\

    \noindent
    Keywords: LEGO Mindstorms EV3, mobile agent, behaviour
\end{abstract}
\end{otherlanguage}


\tableofcontents

\chapter{Wstęp}
\label{ch:wstep}

Celem niniejszej pracy było zaprojektowanie i implementacja aplikacji kontrolującej zachowanie robotów mobilnych. Jest ona odpowiedzialna za nadzór nad zachowaniami robotów (agentów), sterowanie sensorami i efektorami oraz komunikację z wyższymi warstwami architektury, weryfikującymi poprawność całego systemu. Zachowania robota są opisane za pomocą automatów skończonych. Warstwa komunikacyjna dostarcza interfejs przesyłania specjalnych komunikatów za pomocą protokołu UDP.\\

Aplikacja jest zaprojektowana do uruchamiania na systemach Unixowych, w szczególności na klocku centralnym LEGO, ale także na komputerze osobistym w roli nadzorcy. Warstwa połączeniowa jest punktem wspólnym architektury ARM z inną, używając do komunikacji interfejsu gniazd sieciowych.
\clearpage

\section{Motywacja}

Motywacją do podjęcia powyższego tematu było przetestowanie możliwości programowalnych oraz technicznych robota zbudowanego z klocków LEGO Mindstorms EV3. Aplikacja działająca na urządzeniu była napisana z użyciem biblioteki\\
ev3dev \cite{ev3dev} w języku C++.

Możliwość dostępu do sterującego klockiem centralnym systemu Linux zdejmuje ograniczenie używania prostych środowisk graficznych i pozwala osiągnąć dużo więcej małym kosztem. Należało zatem sprawdzić, co najnowsza wersja \\LEGO Mindstorms ma do zaoferowania, w szczególności:
\begin{itemize}
    \item Wydajność napisanych aplikacji z użyciem ww. bilbioteki.
    \item Skuteczność komunikacji z wykorzystaniem bezprzewodowej sieci Wi-Fi.
    \item Dokładność dostarczanych odczytów z sensorów oraz szybkość i niezawodność reakcji na nietypowe zdarzenia.
\end{itemize}

\section{Założenia}

Zostały przyjęte następujące założenia:
\begin{itemize}
  \item Każdy agent jest zdolny do wykonywania pewnych konkretnych zachowań niezależnie od pozostałych agentów.
  \item Zachowania te są reprezentowane za pomocą automatu skończonego.
  \item Dany robot może, ale nie musi być zdolny do wykonania konkretnej czynności. Jest to zależne od podłączonych do niego sensorów i efektorów.
  \item Każdy robot samodzielnie generuje sposób wykonania danej akcji na podstawie dostępnych urządzeń.
  \item Zachowania mogą być dynamicznie tworzone z użyciem specjalnej składni.
  \item Urządzeniem nadzorującym może być inny robot, ale pożądana jest też możliwość kontroli z poziomu zwykłego komputera.
  \item Agenci komunikują się zarówno z jednostką centralną (nadzorująca) jak i między sobą za pomocą sieci bezprzewodowej.
  \item Do poprawnej komunikacji wymagana jest jedna, wspólna dla wszystkich jednostek, istniejąca sieć.
  \item System powinien dostosować się do braków w łączności, umożliwiając komunikację przez pośrednictwo innych agentów.
  \item Agent może poruszać się tylko po płaskiej powierzchni.
\end{itemize}

% \section{Ograniczenia}
%
% \ldots
% TODO

\section{Zawartość rozdziałów}

\begin{description}
  \item[ {~\ref{ch:wstep} \hyperref[ch:wstep]{Wstęp} -} ]zawiera cel pracy razem z założeniami i ograniczeniami.
  \item[ {~\ref{ch:opis_systemu} \hyperref[ch:opis_systemu]{Opis systemu} -} ]opisuje dokładnie aspekty projektowe i technologiczne oraz zawiera wytyczne dotyczące wykorzystanych narzędzi.
  \item[ {~\ref{ch:budowa_aplikacji} \hyperref[ch:budowa_aplikacji]{Budowa aplikacji} -} ]szczegółowo omawia wykorzystane w projekcie klasy, ich wzajemne zależności oraz podział na moduły.
  \item[ {~\ref{ch:zachowania} \hyperref[ch:zachowania]{Zachowania} -} ]zawiera opis budowy oraz działania zachowań złożonych robota wraz z przykładami.
  \item[ {~\ref{ch:komunikacja} \hyperref[ch:komunikacja]{Komunikacja} -} ]opisuje szczegóły implementacji warstwy połączeniowej, bazującej na interfejsie gniazd sieciowych i protokole UDP.
  \item[ {~\ref{ch:testowanie} \hyperref[ch:testowanie]{Testowanie aplikacji} -} ]zawiera przebieg przykładowych testów i ich rezultaty.
  \item[ {~\ref{ch:podsumowanie} \hyperref[ch:podsumowanie]{Podsumowanie} -} ]zawiera wnioski końcowe z przebiegu projektowania aplikacji oraz przeprowadzonych testów.
\end{description}

\chapter{Opis systemu}
\label{ch:opis_systemu}

Przy projektowaniu systemu, należało uwzględnić cechy charakterystyczne wielu różnych dziedzin. Aplikacja łączy w sobie działanie niskopoziomowych, sprzętowych operacji na sensorach i efektorach, zaprogramowaną logikę wyższego poziomu w języku C++ oraz komunikację sieciową w protokole UDP. Architektura ARM oraz docelowa platforma systemowa dodatkowo wpływają na podejmowane decyzje projektowe oraz implementacyjne. Szczegółowy opis użytych technologii znajduje się poniżej.

\section{Wymagania}
\begin{itemize}
    \item Użyte oprogramowanie działa na klocku centralnym LEGO Mindstorms EV3
    \item Wybrana do projektu biblioteka umożliwia pracę na urządzeniach podłączanych do klocka centralnego przy użyciu języka wyższego poziomu
    \item Aplikacja umożliwia konstrukcję zachowań złożonych w postaci automatu skończonego
    \item Zachowania robota są parametryzowane
    \item Robot powinien reagować na generowane zdarzenia w jak najkrótszym czasie
    \item Robot poprawnie reaguje na odłączenie, bądź zaburzenie pracy któregokolwiek z podłączonych urządzeń
    \item Oprogramowanie monitoruje poziom baterii i reaguje na zmiany jego poziomu
    \item Oprogramowanie umożliwia komunikację z innymi robotami za pomocą sieci bezprzewodowej
    \item Komunikacja jest szybka i odporna na błędy oraz gubienie pakietów
    \item Kod aplikacji jest czytelny oraz utrzymuje niezmienną konwencję\\nazewnictwa
    \item Kod aplikacji, w szczególności pliki nagłówkowe, jest dobrze\\udokumentowany.
\end{itemize}

\section{Użyta technologia}

\subsection{Dostępne środowiska}

Istnieje wiele dostępnych środowisk na licencji otwartego oprogramowania, dedykowanych LEGO Mindstorms EV3, z których każde ma trochę inne zastosowanie. Poniżej opisane są najbardziej popularne z nich.

\subsubsection{leJOS}

Projekt leJOS \cite{lejos} jest oprogramowaniem zastępującym domyślnie zainstalowany system na klocku centralnym. Bazuje on na maszynie wirtualnej języka Java i jest kompatybilny z platformami Linux, Windows oraz Mac OS X. Programowanie aplikacji oraz transfer plików wykonywalnych do robota odbywa się za pomocą wtyczki do środowiska Eclipse lub za pomocą linii poleceń.

Środowisko udostępnia wiele możliwości języka Java, w szczególności obsługę synchronizacji i wyjątków, wątki oraz większość klas z pakietów java.lang, java.util oraz java.io. Dostarcza ono także ujednolicony interfejs do obsługi sensorów. Na\,wyświetlaczu klocka centralnego udostępniony jest nowy interfejs graficzny z wieloma dodatkowymi funkcjonalnościami. Ponadto, leJOS posiada wiele zdefiniowanych modułów obsługujących połączenia z wieloma robotami, algorytmy lokalizacyjne czy dostęp do konsoli klocka LEGO. Co ważniejsze, są one uruchamiane na\,zwykłym komputerze, a nadzór odbywa się zdalnie.

Mimo oferowanego bogactwa możliwości, środowisko leJOS nie byłoby dobrym wyborem. Zastąpienie oryginalnego oprogramowania klocka, dodatkowe wykorzystanie zasobów urządzenia przez maszynę wirtualną, czy brak niskopoziomowego wsparcia to główne powody dla których ta biblioteka została odrzucona.

\subsubsection{MonoBrick}

Innym wieloplatformowym środowiskiem, które współpracuje z zestawem Mindstorms EV3, jest biblioteka .NET o nazwie MonoBrick \cite{monobrick}. Wspiera ona języki takie jak C\#, F\# czy IronPython. Instalacja oprogramowania nie zastępuje domyślnego systemu na klocku, ale uruchamiana jest oddzielnie z poziomu karty SD. Oferowane możliwości są prawie tak duże, jak w przypadku leJOS, aczkolwiek projektowi brakowało solidnej dokumentacji, a także wiele odnośników na stronie prowadziło do nieistniejących witryn.

Powoływanie się na niepewną bibliotekę, która nie umożliwia natywnej kompilacji kodu, byłoby dużym błędem. Projekt prężnie się rozwija, ale w trakcie wyboru docelowego oprogramowania nie był wystarczająco stabilny i funkcjonalny.

\subsubsection{ev3dev}

Najbardziej dogodnym rozwiązaniem okazał się projekt ev3dev. Jest to dopasowana do potrzeb klocka LEGO dystrybucja Linuxa (Debian Jessie), która jest wgrywana na kartę SD i uruchamiana obok istniejącego systemu. Zawsze istnieje możliwość przywrócenia domyślnego stanu klocka przez wyjęcie karty z systemem. Platforma stworzona w ramach ev3dev zawiera wiele sterowników, nie tylko do\,akcesoriów zestawu EV3, ale także poprzednich dystrybucji LEGO Mindstorms oraz komponentów wytwarzanych przez osoby trzecie. Możliwe jest programowanie klocka w języku C/C++, ale ev3dev obsługuje też wiele innych języków, takich jak Python czy Lua. To\,wszystko daje dużą swobodę samego programowania, jak\,i sposobu tworzenia programu i komunikacji z urządzeniem. Kompilacja aplikacji może odbywać się bezpośrednio na urządzeniu lub na komputerze z wbudowanym kompilatorem na procesory typu ARM. Komunikacja z klockiem centralnym realizowalna jest na trzy sposoby: za pomocą Wi-Fi, Bluetooth lub przy użyciu kabla USB.

Zalety, takie jak praca w natywnym, obiektowym języku na natywnej platformie, niskopoziomowość, pełna kontrola nad sprzętem oraz brak narzutów maszyn wirtualnych zadecydowały o przyjęciu ev3dev jako biblioteki wykorzystanej w tym projekcie.

\subsection{Sprzęt}

Parametry klocka centralnego:
\begin{description}
    \item[System:] Debian Jessie
    \item[Wyświetlacz:] 178x128 pikseli, monochromatyczny
    \item[Procesor:] TI Sitara AM1808 300MHz
    \item[Pamięć:] 64MB RAM, 16MB Flash, slot na kartę microSDHC
    \item[Sieć:] WiFi, Bluetooth, USB.
\end{description}

\subsection{Konfiguracja}

Każda jednostka z wgraną biblioteką ev3dev posiada użytkownika {\tt robot} z hasłem {\tt maker}. Konfiguracja nowego systemu odbyła się w kilku krokach:

\begin{enumerate}
    \item Na kartę microSDHC wgrany został specjalne spreparowany obraz systemu, pobrany ze strony głównej projektu ev3dev
    \item Pierwsze uruchomienie robota z wgranym obrazem na karcie było poprzedzone długim przygotowaniem wszystkich plików
    \item Przy użyciu połączenia SSH przez kabel USB, system został skonfigurowany i pobrane zostały wszystkie wymagane pakiety
    \item Dalsza komunikacja odbywała się bezprzewodowo z użyciem urządzenia NETGEAR WNA1100, podłączonego do portu USB klocka centralnego
    \item Aplikacja była kompilowana na laptopie z systemem Ubuntu 15.10 i synchronizowana zdalnie z robotem.
\end{enumerate}

\subsection{Biblioteka}
W ramach projektu ev3dev dostępne są dwa pliki źródłowe napisane w języku C++. Dostarczają one wymagany interfejs do sterowania klockiem centralnym i\,podłączonymi do niego urządzeniami.\\

\noindent Wersja użytej biblioteki: 0.9.3-pre, rev 2.

\subsubsection{Wprowadzone zmiany}
Początkowo, w wersji 0.9.2-pre, biblioteka nie była kompatybilna ze wszystkimi urządzeniami dostarczonymi przez LEGO. Wymagane było dopisanie rozpoznawalnych nazw sterowników dla sensorów oraz zaimplementowanie własnej obsługi diod LED przedniego panelu, w szczególności funkcji migania. Ponadto dopisane zostało specjalne pole klasy {\tt ev3dev::motor} pamiętające jego typ.

W finalnej fazie projektowania, biblioteka została zaktualizowana do najnowszej wersji 0.9.3-pre, co zdjęło konieczność posiadania specjalnych nazw sterowników, a także znacząco przyspieszyło obsługę interfejsu klocka oraz transmisję danych.

\subsection{Możliwości}
Użyte środowisko Linux orax język programowania C++ dostarczają wiele możliwości programistycznych, a w szczególności:

\begin{itemize}
    \item Użycie biblioteki {\tt stl} i zgodność ze standardem C++11
    \item Wątki
    \item Polimorfizm
    \item Komunikację przez protokół UDP z wykorzystaniem gniazd.
\end{itemize}

\subsection{Narzędzia}
Projekt aplikacji był rozwijany z wykorzystaniem narzędzie NetBeans. Mimo iż sama aplikacja może być skompilowana z poziomu konsoli i narzędzia Makefile, NetBeans dostarcza także wygodne narzędzia debugujące oraz przyspiesza pisanie kodu.

\subsubsection{Konfiguracje aplikacji}
Zostały zdefiniowane dwie domyślne konfiguracje:
\begin{itemize}
    \item D\_ARM: konfiguracja przeznaczona na urządzenia z procesorami typu ARM. Domyślnie przeznaczona do uruchamiania na robocie mobilnym
    \begin{description}
        \item[{\small Kompilator:}] {\small arm-linux-gnueabi-g++}
        \item[{\small Flagi kompilacji:}] {\small -D\_GLIBCXX\_USE\_NANOSLEEP -pthread -static-libstdc++ -std=c++11 -DAGENT}
    \end{description}
    \item D\_DESKTOP: konfiguracja kompilowana z myślą o tradycyjnych komputerach osobistych. Domyślnie przeznaczona do uruchomienia w trybie nadzorcy systemu, komunikującego się zdalnie z robotami.
    \begin{description}
        \item[{\small Kompilator:}] {\small g++}
        \item[{\small Flagi kompilacji:}] {\small -D\_GLIBCXX\_USE\_NANOSLEEP -pthread -static-libstdc++ -std=c++11}
    \end{description}
\end{itemize}
\noindent Obie konfiguracje posiadają dodatkowo wersję z przedrostkiem {\tt R\_}, które oznaczają wersję Release zamiast wersji Debug.\\

Inne użyte narzędzia to przede wszystkim aplikacja Doxygen służąca generowaniu dokumentacji kodu programu, system kontroli wersji \gls{git} do zarządzania całym projektem oraz oprogramowanie LaTeX, za pomocą którego wygenerowany został ten dokument.

\subsection{Konstrukcja robota}
W celu przetestowania zaimplementowanych funkcjonalności, wyposażono robota w następujące elementy:
\begin{itemize}
    \item Dwa duże motory do poruszania się po płaskiej powierzchni
    \item Ultradźwiękowy sensor odległości ustawiony w kierunku poruszania się
    \item Przedni zderzak oraz sensor dotyku do wykrywania zderzeń
    \item Sensor koloru do wykrywania zmian w odcieniu powierzchni.
\end{itemize}

\chapter{Budowa aplikacji}

\section{Moduły i klasy}
\indent \indent W celu uzyskania przejrzystości aplikacji, wydzielone zostały moduły\footnote{W kontekście tej pracy \textit{moduł} oznacza pewną grupę skojarzonych ze sobą klas.}, które opisują pewien fragment funkcjonalności programu. Moduły niższych warstw mogą być wykorzystywane przez moduły warstw wyższych lub być tylko zestawem dodatkowych narzędzi. Kolejne punkty opisują w czym dany moduł się specjalizuje i jakie klasy wchodzą w jego skład. Szczegółowe dane na temat klas oraz ich metod i pól znajdują się w dokumentacji kodu.
% TODO Jakoś podlinkować czy wspomnieć o tej docu, może gdzieś na końcu.

\subsection{Komendy}
\indent \indent Klasy komend są tak naprawdę nakładką na istniejące mechanizmy biblioteki ev3dev, operujące bezpośrednio na sprzęcie. Oprócz właściwej komendy, będącej poleceniem dla efektora bądź sensora, dana klasa zawiera referencje do obiektu, na którym ma zostać wykonana oraz jej parametry, o ile takowe posiada. Nazewnictwo klas dokładnie odwzorowuje nazwy komend przekazywanej urządzeniom. W obrębie konkretnych komend, definiowane są także stałe opisujące charakter przekazywanych argumentów oraz ich limity.\\
Komendy zostały podzielone na dwie podgrupy:
\begin{description}
    \item[Komendy motorów:] Klasa bazowa - {\tt CommandMotor}. Zawierają referencje do klasy {\tt Motor} oraz opcjonalnie przechowują także przekazywane parametry. \\Np przykład: {\tt CommandMotorStop}, {\tt CommandMotorRunForever}.

    \item[Komendy sensorów:] Klasa bazowa - {\tt CommandSensor}. Zawierają referencje do klasy {\tt Sensor}. Definiują obsługiwane tryby danego sensora. Komendy te nie służą do pobierania wartości, lecz tylko do zmiany ich ustawień. Pobieranie wartości używane jest przy pomocy specjalnej klasy {\tt Devices}. \\Przykładowa komenda: {\tt CommandSensorSetMode}.
\end{description}
Klasą bazową dla wszystkich komend jest {\tt Command}.

\subsection{Akcje}
\indent \indent Akcje są kolejnym stopniem abstrakcji definiowania zachowań robota. Klasy akcji przechowują przede wszystkim sekwencje komend, które mają zostać wykonane. Ponadto, z powodu natychmiastowego charakteru wykonywania wszystkich zgromadzonych komend, akcja może mieć zdefiniowany warunek jej zakończenia. Przyjmuje ona postać funkcji anonimowej, w której następuje zwrócenie wartości prawda lub fałsz na podstawie dowolnie sprecyzowanych instrukcji. Pozwala to wyższej warstwie sterującej sprawdzić, czy kolejna akcja może zostać wykonana. Dodatkowo, akcje mogą deklarować dopuszczalne zdarzenia, które przerywają jej działanie lub zmieniają jej parametry.

Wszystkie dostępne klasy akcji są zdefiniowane w aplikacji i nie istnieje możliwość zwiększenia zbioru o nowe bądź dynamicznego generowania nowych, własnych klas. Ta decyzja implementacyjna jest podyktowana specyfiką konkretnych modeli robotów, różniących się budową oraz podłączonymi akcesoriami, a co za tym idzie, odmiennym sposobem implementacji tych samych czynności.

Klasy opisujące konkretne akcje, np. {\tt ActionDriveDistance}, dziedziczą po klasie {\tt Action}. Wspólnymi elementami każdej z nich są: typ, warunek końcowy, sekwencja komend oraz metody wykonawcze. Różnią się natomiast dodatkowymi parametrami, takimi jak prędkość czy kąt obrotu. Ponadto, istnieje możliwość wygodnego zapętlenia jednej lub wielu akcji dowolną liczbę razy za pomocą specjalnej klasy - {\tt ActionRepeat}. Konstruktor tej klasy przyjmuje liczbę powtórzeń oraz listę akcji, które zostaną wykonane w podanej kolejności.

\subsection{Zachowania}

\indent \indent Definiowanie zachowań jest dużo trudniejsze niż akcji czy komend, gdyż bazują one na schemacie automatu skończonego. Oprócz konkretnych akcji, na jakich dane zachowanie ma się opierać, należy zdefiniować przejścia pomiędzy stanami (akcjami) w toku poprawnego wykonania oraz specjalne warunki zmiany stanu w reakcji na zaistniałe zdarzenia (np. napotkana przeszkoda lub utrata połączenia).

Podobnie jak w przypadku akcji i komend, każde zachowanie zdefiniowane jest w osobnej klasie (np. {\tt BehaviourExplore}), które dziedziczy po wspólnej klasie {\tt Behaviour}. Każde z nich zawiera typ, strukturę stanów automatu złożonych z akcji i przejść oraz funkcje wykonawcze. Ponadto zachowania pochodne zawierają własne parametry, np. maksymalny dystans do przejechania.

% TODO Diagram przykładowego zachowania.

\subsection{Urządzenia}

Biblioteka ev3dev dostarcza wygodnego interfejsu do zarządzania urządzeniami przez wygenerowanie drzew klas dla sensorów i efektorów. W ramach tej aplikacji, na każdy typ urządzenia została nałożona specjalna klasa pośrednicząca, w środku której dopiero znajduje się referencja do właściwego obiektu. Są to klasy {\tt Motor} oraz {\tt Sensor}, które po odpowiedniej identyfikacji zostają zmapowane do par port-urządzenie.

Całością nadzoruje klasa {\tt Devices} napisana zgodnie ze wzorcem projektowym Singleton. Ograniczone są w ten sposób nadmierne kopie obiektów i potencjalnie niejednoznaczne odwołania. Ponadto klasy urządzeń są potrzebne w wielu różnych miejscach aplikacji, a wzorzec ten umożliwia taki dostęp za pomocą statycznego wydobycia instancji.

Moduł urządzeń jest także odpowiedzialny za detekcję zdarzeń. Wyższe warstwy mogą zgłaszać zdarzenia, na które klasa {\tt Devices} ma nasłuchiwać. Jeśli dane zdarzenie wystąpi, wysyłane jest do odpowiedniej kolejki dla wyższych warstw do przetworzenia. Zgłaszane mogą być również zdarzenie niezależne, np. niski poziom baterii lub odłączenie urządzenia.

\subsection{Robot}

Jest to najbardziej rozbudowana klasa, ponieważ agreguje w sobie działanie wielu modułów. Zarządza zarówno urządzeniami podłączonymi do klocka centralnego, steruje zachowaniem robota oraz przetwarza przesłane komunikaty oraz zdarzenia. Główna metoda klasy {\tt run} jest uruchamiana w głównym wątku aplikacji i w pętli przetwarza wszystkie dane. Jest bezpośrednio zsynchronizowana z wątkiem komunikacyjnym za pomocą dwóch kolejek wiadomości - nadawczej i odbiorczej. Wyróżnienie dwóch niezależnych ścieżek wykonania umożliwia lepsze zarządzanie zasobami sprzętowymi oraz pozwala robotowi na samodzielne działanie, niezależnie od przesyłanych pakietów.

Robot może być fizycznie zbudowany na wiele różnych sposobów. Dlatego wymagane jest, żeby zdefiniowane były konkretne klasy implementujące szczegóły danego modelu. W związku z powyższym, klasa bazowa {\tt Robot} jest klasą abstrakcyjną, a każdemu wariantowi konstrukcji odpowiada osobna klasa podrzędna. W celu ujednolicenia interfejsu, każdy model deklaruje listę obsługiwanych akcji oraz wymaganych do tego podłączonych urządzeń. W ten sposób można łatwo sprawdzić, czy dane zachowanie jest obsługiwane i jakich pomiarów robot może dostarczyć. Klasy konkretnych modeli, bazując na zdefiniowanych funkcjach wirtualnych, definiują swoje wersje w sposób adekwatny do domniemanego efektu końcowego wymaganych akcji.

Każdy instancja robota posiada również maszynę stanów oraz zdefiniowane warunki przejść pomiędzy nimi. Każdy stan przetwarza tylko konkretne komunikaty i zdarzenia. Rezultatem ich przetworzenia może być wysłanie specjalnej wiadomości zwrotnej lub zmiana stanu na inny. To ostanie wiąże się jeszcze ze zmianą koloru oraz częstotliwości migania diod LED umieszczonych na przednim panelu klocka centralnego.

% TODO Diagram stanów + kolory

\subsection{Komunikacja}

Komunikacja pomiędzy różnymi partiami całego systemu odbywa się na dwóch poziomach:

\subsubsection{Zdarzenia}

Poszczególne moduły robota, takie jak akcje czy klasa urządzeń, mogą generować zdarzenia. Kolejka zdarzeń, będąca Singletonem,...
> wiele punktow wejscia
> 1 punkt wyjscia
> klasa Event
> takze wyjatki

\subsubsection{Komunikacja UDP}

\subsection{Nadzorca}

\subsection{Moduły dodatkowe}

\chapter{Zachowania}
\label{ch:zachowania}

Zachowania definiują najwyższy stopień abstrakcji sterowania agentem. Oprócz właściwych akcji, które wykonują, mogą mieć zdefiniowane również zdarzenia, które wymuszają niestandardowy przebieg wykonania. Ponadto, zlecają również śledzenie konkretnych wartości na podłączonych urządzeniach.

\section{Budowa zachowań złożonych}

Każde zachowanie składa się z kilku elementów:
\begin{itemize}
    \item Listę akcji, z których każda stanowi de facto osobny stan.
    \item Listę możliwych przejść między stanami wraz z warunkami.
    \item Zdarzenia, na które zachowanie będzie reagowało w określony sposób.
\end{itemize}

\subsection{Tworzenie zachowań}

Stworzenie zachowania może odbyć się na dwa sposoby. Pierwszym z nich jest skorzystanie z zachowania już zdefiniowanego w bibliotece. W tym przypadku zbierane są ustalone wcześniej, odpowiednio ukonkretnione akcje razem z warunkami przejść. Drugim wariantem jest stworzenie własnego zachowania. Jeśli taki typ zostanie wykryty, do agenta należy dostarczyć dodatkowe dane opisujące użyte akcje, przejścia i zdarzenia.

% TODO XML ?

\subsection{Sterowanie wykonywanymi akcjami}

\subsection{Przejścia pomiędzy stanami}


\section{Wykonywanie zachowania}

% TODO Delete those subsections if not enough text
\subsection{Wykonywanie akcji}
\subsection{Reakcja na zdarzenia}
\subsection{Zbieranie danych}


\section{Przykładowe zachowania}

\subsection{Zwiedzanie otoczenia}

\chapter{Komunikacja}
\label{ch:komunikacja}

\section{Komunikacja przez sieć}

Komunikacja przez sieć odbywa się przy użyciu protokołu UDP oraz z wykorzystaniem interfejsu gniazd systemu Unix. Domyślnie, ustawione są następujące parametry:
\begin{itemize}
    \item Numer portu: 12345
    \item Maksymalny rozmiar paczki: 4kB
    \item Liczba wysyłanych kopii pakietów: 3
    \item Wielkość kolejki przechowującej ostatnio odebrane pakiety: 50
\end{itemize}

W klasie {\tt CommUtils} zdefiniowane są metody {\tt sendMessage, receiveMessage} oraz {\tt receiveMessageDelay}. Przetwarzają one zdefiniowane struktury aplikacji na niskopoziomowe bufory danych wysyłane lub odbierane za pomocą systemowych funkcji (Rysunek \ref{fig:low_level_comm}). W trakcie ich wykonywania sprawdzane są dostępne interfejsy sieciowe oraz następuje tworzenie gniazd na zadanym porcie. Metoda {\tt receiveMessage} jest blokująca, natomiast {\tt receiveMessageDelay} czeka na odpowiedź tylko przez zadany w milisekundach interwał.

\begin{figure}[!ht]
    \centering
        \includegraphics[width=1\textwidth]{diagrams/low_level_comm.png}
    \caption{Uproszczony diagram niskopoziomowej komunikacji.\label{fig:low_level_comm}}
\end{figure}

\subsection{Synchronizacja z nadzorcą}

Na samym początku, tuż po uruchomieniu agentów i nadzorcy, następuje synchronizacja. Nadzorca nasłuchuje sieć, oczekując na wiadomości typu AGENT. Uruchomiony agent jest w stanie IDLE i wysyła wiadomość typu AGENT na adres rozgłoszeniowy z losowym identyfikatorem agenta, oczekując na komunikat MASTER. Czekanie nie trwa dłużej niż zadany interwał czasu (domyślnie nadzorca ma 10 sekund na odpowiedź). Nadzorca odbiera żądanie od agenta i dodaje go do swojej bazy, jednocześnie przypisując mu identyfikator oraz zapamiętując adres ip. Następnie na adres wysyłany jest komunikat MASTER z nowo przydzielonym identyfikatorem wpisanym w miejsce odbiorcy. Jeśli agent przyjmie taki pakiet, zwraca wiadomość typu ACK. Po takiej wymianie wiadomości (rysunek \ref{fig:comm_synch}), dany robot jest zsynchronizowany, ma przypisany numer identyfikujący oraz jest odnotowany w bazie. Następuje przejście do stanu ACTIVE, oznaczające danego agenta, jako aktywnego w systemie.

\begin{figure}[!ht]
    \centering
        \includegraphics[width=0.75\textwidth]{diagrams/comm_synch.png}
    \caption{Diagram komunikacji dla poprawnego rozpoczęcia zachowania.\label{fig:comm_synch}}
\end{figure}

W przypadku braku aktywnego agenta, nadzorca pozostaje bezczynny. Jeśli\,aktywny agent nie zostanie zsynchronizowany, przechodzi on do stanu PANIC (rysunek \ref{fig:comm_idle_panic}, patrz:~\ref{sec:utrata_polaczenia} \hyperref[sec:utrata_polaczenia]{Utrata połączenia z nadzorcą}).

\begin{figure}[!ht]
    \centering
        \includegraphics[width=0.75\textwidth]{diagrams/comm_idle_panic.png}
    \caption{Diagram komunikacji dla braku synchronizacji z nadzorcą.\label{fig:comm_idle_panic}}
\end{figure}

W każdym stanie oprócz stanu PANIC, z zadanym interwałem czasowym wysyłane są do nadzorcy wiadomości podtrzymujące połączenie. Agent wysyła komunikat PING, nadzorca natychmiast po jego odebraniu odsyła pakiet PONG. Jeśli oczekiwanie przekroczy zadany próg czasowy, następuje przejście do nowego stanu (rysunek {\ref{fig:comm_active_panic}}).

\begin{figure}[!ht]
    \centering
        \includegraphics[width=0.75\textwidth]{diagrams/comm_active_panic.png}
    \caption{Diagram komunikacji dla utraty połączenia po synchronizacji.\label{fig:comm_active_panic}}
\end{figure}

\subsection{Oczekiwanie na aktywację}

W stanie ACTIVE agent oczekuje na przydzielenie zadania. Jeśli takie nie nadchodzi, wysyłane są pakiety podtrzymujące połączenie. W przypadku braku komunikatów zwrotnych, robot przechodzi do stanu PANIC (rysunek \ref{fig:comm_active_panic}).

Nadzorca decyduje o przydzieleniu konkretnego zachowania danemu agentowi. Wysyła do niego wiadomość typu BEHAVIOUR z parametrami zawierającymi jego typ oraz ewentualne dodatkowe argumenty. Jeśli agent zdolny jest do wykonywania danego zadania, zwraca pakiet ACK, w przeciwnym razie pakiet NOT. Komunikat ACK skutkuje odpowiedzią START, po której agent przechodzi do stanu WORKING. W przypadku komnunikatu NOT, nadzorca może zdefiniować inne zadanie lub pozostować agenta tymczasowo nieaktywnego (rysunek \ref{fig:comm_synch}).

\subsection{Przetwarzanie zachowania}

Agent może przetwarzać zlecone instrukcje tylko w stanie WORKING. Po każdej prawidłowo zakończonej akcji, wysyłana jest wiadomość typu ACTION\_OK z identyfikatorem poprawnie wykonanej akcji. Pozwala to nadzorcy logować przebieg pray agentów. W przypadku zakłócenia przebiegu wykonania, robot może wysłać pakiet ACTION\_INTERR, jeśli zaszło jakieś zdarzenie wymagające zmiany postępowania, a także pakiet ACTION\_ERR, jeśli wystąpił błąd. W pierwszym przypadku, jeśli robot sam nie może określić następnej czynności, następuje oczekiwanie na pakiet CONTINUE. W drugim przypadku, robot bezwarunkowo zawiesza aktualną działalność i oczekuje na wiadomość RESTART, która rozpocznie przebieg bieżącego zachowania od nowa lub w przypadku wiadomości BEHAVIOUR, przebieg nowego zachowania. Niezależnie od rodzaju wiadomości wznawiającej, agent odpowiada komunikatami ACK lub NOT.

W czasie wykonywania akcji, pobierane są odczyty z sensorów. Jeśli zachowanie ma zdefiniowane konkretne sensory, których pomiary są potrzebne, wysyłane są wiadomości typu SENSOR\_VALUE wraz z pomiarem (rysunek \ref{fig:comm_working_process}). Nadzorca interpretuje je i w miarę możliwości zapamiętuje. Mogą one zadecydować o zmianie zachowania danego robota.

\begin{figure}[!ht]
    \centering
        \includegraphics[width=0.75\textwidth]{diagrams/comm_working_process.png}
    \caption{Przykładowy przebieg dla stanu WORKING.\label{fig:comm_working_process}}
\end{figure}

Brak odpowiedzi na wiadomości podtrzymujące lub wysłanie wiadomości typu PAUSE skutkuje przejściem do stanu PAUSED (rysunek \ref{fig:comm_working_panic}. Otrzymanie pakietu STAND\_BY powoduje przejście agenta z powrotem w stan ACTIVE.

\begin{figure}[!ht]
    \centering
        \includegraphics[width=0.75\textwidth]{diagrams/comm_working_panic.png}
    \caption{Brak komunikacji w stanie WORKING.\label{fig:comm_working_panic}}
\end{figure}

\subsection{Bezczynność agenta}

W stanie PAUSED, robot pamięta aktualny przebieg zachowania ale nie wykonuje żadnych czynności. Dopiero po otrzymaniu komunikatu RESUME i odesłaniu odpowiedzi ACK, może powrócić do stanu WORKING i kontynuować wstrzymane zachowanie.

Brak wiadomości podtrzymujących od strony nadzorcy powoduje przejście do\,stanu PANIC (rysunek \ref{fig:comm_pause_process}).

\begin{figure}[!ht]
    \centering
        \includegraphics[width=0.75\textwidth]{diagrams/comm_pause_process.png}
    \caption{Diagram komunikacji dla stany pauzy.\label{fig:comm_pause_process}}
\end{figure}

\section{Komunikacja wewnątrz jednostki}

W obrębie jednego agenta, wyróżnia się dwie podstawowe formy komunikacji: wiadomości i zdarzenia.

\subsection{Wiadomości}

Przy uruchomieniu programu, tworzony jest dodatkowy wątek komunikacyjny, który jest zsynchronizowany z wątkiem głównym robota (lub nadzorcy). Wymiana informacji odbywa się za pomocą dwóch, przeciwnie przydzielonych obiektów klasy szablonowej {\tt Queue}, która zawiera interfejs obsługi kolejki wiadomości wraz z operacjami jednoczesnego dostępu wielu wątków. Kolejka wiadomości odbieranych przez klasę {\tt Robot} jest kolejką wiadomości wysyłanych przez klasę {\tt Communication} i odwrotnie w przypadku drugiego obiektu.

Po stronie komunikacji odbywa się jedynie przesyłanie i odbieranie pakietów między fizycznymi urządzeniami oraz sprawdzanie duplikatów. Interpretacją zajmuje się wątek główny. Wszystkie możliwe do wysłania pakiety są zapisywane w\,klasie {\tt Message}, w której skład wchodzą następujące elementy:
\begin{itemize}
    \item Identyfikator nadawcy, przydzielany automatycznie
    \item Identyfikator odbiorcy, stały i równy 1 w przypadku, gdy odbiorcą jest nadzorca. W przeciwnym razie wybierany spośród agentów z bazy
    \item Identyfikator wiadomości, sukcesywnie inkrementowany
    \item Rodzaj wiadomości, bedący wartością typu wyliczeniowego, czyli tak naprawdę liczbą
    \item (opcjonalnie) Dodatkowe parametry w zależniości od typu wiadomości.
\end{itemize}

Przed wysłaniem, wiadomość będąca obiektem zamieniana jest na łańcuch znaków zawierający jej prototyp. Przykładowy prototyp może wyglądać następująco:
\begin{center} {\tt 14:1:254:5:foo:bar},\\ \end{center}
gdzie wartości oddzielone separatorem (domyślnie dwukropek) oznaczają kolejno: identyfikator agenta, identyfikator odbiorcy (nadzorcy), numer wiadomości, typ\,wiadomości oraz dodatkowe parametry.

Przy odbieraniu kolejnych wiadomości, w buforze cyklicznym zapisywane są ich prototypy. Nowe wiadomości zastępują stare, jeśli bufor jest w całości zapełniony. Przed dodaniem prototypu do bufora, następuje sprawdzenie, czy nie jest on już w\,nim zapisany. Jeśli tak, oznacza to odebranie duplikatu i odrzucenie wiadomości. Rozwiązanie to jest konieczne w przypadku wielu agentów, gdyż kopie wysyłanych przez nich wiadomości mogą przychodzić naprzemiennie.

Odebrany i zakwalifikowany do przyjęcia prototyp zostaje zdekodowany do postaci obiektu klasy {\tt Message} i przesłany do wątku głównego w celu interpretacji. Pakiety, które zawierają niezaktualizowane identyfikatory wiadomości lub błędne informacje o nadawcy są odrzucane. Ma to miejsce szczególnie w przypadku pakietów wysyłanych na adres rozgłoszeniowy, np. podczas synchronizacji agenta z\,nadzorcą.

\subsection{Zdarzenia}

W celu usprawnienia komunikacji oraz zwiększenia czytelności całej aplikacji, konkretne warstwy i moduły komunikują się z klasą {\tt Robot} za pomocą zdarzeń opisanych klasą {\tt Event}. Kolejka przechowująca zdarzenia jest wspólna dla całej instancji aplikacji. Dodawać nowe elementy mogą klasy zachowań, akcji czy urządzeń, a jedynym punktem wyjścia jest główna klasa robota (rysunek \ref{fig:event_queue}), który interpretuje zdarzenie i reaguje na nie.

Rozróżniane są zdarzenia zwykłe oraz krytyczne, które przerywają wykonywanie części bądź wszystkich działań. Do takich zdarzeń należą np. odłączenie urządzenia czy niski poziom baterii robota. Pozostałe zdarzenia mogą przyjmować formę czysto informacyjną (np. poprawnie podłączone urządzenia), ostrzegającą (np. wykryto zderzenie) lub zgłaszającą wystąpienie wyjątku. Taka forma wewnętrznej komunikacji zdejmuje konieczność przekazywania zdarzeń w argumentach wielu funkcji lub pamiętania bezpośrednich referencji do obiektów generujących zdarzenia oraz pozwala łatwo koordynować cały cykl życia obiektu {\tt Event} z poziomu jednej kolejki. Ponadto, dodanie lub usunięcie nowego obiektu generującego bądź nasłuchującego jest proste i nie ingeruje zbyt mocno w istniejącą strukturę komunikacji.

% TODO Diagram przykładowych zdarzeń

\section{Utrata połączenia z nadzorcą}
\label{sec:utrata_polaczenia}

Zdarzenie to może zajść w dwóch wariantach:

\subsection{Jeden agent traci połączenie z nadzorcą}

Robot przechodzi do stanu PANIC i przez adres rozgłoszeniowy informuje pozostałych agentów o utraconym połączeniu. Jeden z nich może odpowiedzieć pozytywnie, wysyłając potwierdzenie posiadania komunikacji z nadzorcą. Wtedy następuje synchronizacja dwóch agentów i ustalenie przekierowania wiadomości przez robota pośredniego (Rysunek ...). Jeżeli po pewnym czasie, nikt nie odpowie pozytywnie oraz nie zacznie się głosowanie opisane w punkcie poniżej, robot kończy swoje działanie.

% TODO Rysunek

\subsection{Wszyscy agenci tracą połączenie z nadzorcą}

Wszystkie roboty utraciły połączenie i żaden nie otrzymał odpowiedzi na zapytanie o przekierowanie. W takim przypadku odbywa się głosowanie. Każdy agent wysyła na adres rozgłoszeniowy swój osobisty wynik\footnote{Wynik agenta zależy m. in. od ilości przesłanych danych, poprawności wykonanych akcji czy czasu aktywności.}. Po pewnym czasie, kiedy wszystkie roboty znają już swoje wyniki, ten z najniższym rezultatem mianuje siebie nadzorcą. W przypadku tych samych wyników, o wyborze decyduje czynnik losowy. Wybierany zostaje najsłabszy rezultat, ponieważ warto pozwolić dobremu agentowi na kontynuację działań.

Problem jaki może napotkać grupa robotów, to jednoczesny wybór dwóch nadzorców. Ograniczenie użycia jednej wspólnej sieci nie eliminuje problemu gubienia pakietów, a przez to potencjalnej sytuacji, w której mamy wiele niezależnych głosowań, bądź głosowania są niepełne. W tym celu, oprócz wysłania przez danego agenta swoich danych, wysyła on dodatkowo liczbę oznaczającą ilość otrzymanych rekordów. W ten sposób agent A może dowiedzieć się od agenta B, że nie dostał informacji od innego agenta, np. C, kiedy połączenie między A i C jest niemożliwe. Ponadto, własny wynik jest ponownie rozsyłany za każdym zwiększeniem ilości informacji o pozostałych agentach. Jednostki nie posiadające kompletu informacji nie biorą udziału w finalnym głosowaniu, ponieważ mogą podjąć złą decyzję, a ponadto utrudniona komunikacja jest dla nich dodatkowo niekorzystna. Cały schemat działania opisuje Rysunek ... .

% TODO rysunek

\chapter{Testowanie aplikacji}
\label{ch:testowanie}

\section{Testy poprawności}

\subsection{Testowanie pojedynczych akcji}

\begin{description}
    \item[Ruch z zadaną prędkością -]wykonanie akcji powoduje utrzymywanie stałej prędkości obu motorów przez czas nieograniczony
    \item[Zatrzymywanie -]robot wprawiony w ruch zatrzymuje się niemal natychmiastowo. W zależności od sposobu zatrzymywania się, czas po którym prędkość jest równa zeru, a co za tym idzie, dystans hamowania, zmienia się. W\,przypadku trybu {\tt coast} jest on największy i równy od 1 do 5 cm w zależności od szybkości jazdy. Zaś w przypadku trybu {\tt hold}, koła zatrzymują się natychmiast i dystans maleje do maksymalnie 2 cm na śliskiej powierzchni
    \item[Ruch na zadaną odległość -]akcja wykonywana poprawnie z dokładnością do\,kilku nadmiarowych centymetrów. Nadwyżka odległości zależy od ustalonego trybu zatrzymywania oraz rodzaju nawierzchni, a wynika z opóźnienia odczytów z efektorów. Na podstawie zmierzonego promienia koła, zadanej odległości oraz ilości jednostek w jednym pełnym obrocie, można wyliczyć docelową pozycję motoru i zatrzymać go, gdy ta zostanie osiągnięta
    \item[Obrót -]jest wykonywany wokół osi równo oddalonej od obu kół. Podobnie jak w przypadku ruchu, sprawdzana jest zmiana położenia motoru, której docelowa wartość wyliczana jest na podstawie rozstawu kół i kąta. Gdy jeden efektor ma zadaną wartość dodatnią prędkości, drugi ma wartość przeciwną. Wykryto odchylenie o kilka stopni przy większych prędkościach. Zalecana jest zatem mała prędkość obrotu, w celu uzyskania większej precyzji.
\end{description}

\subsection{Testowanie obsługi zdarzeń}

\subsubsection{\underline{Rozpoczęcie przetwarzania zachowania}}


\begin{description}
    \item[Warunki:]Robot implementuje wymagane przez zachowanie akcje, posiada urządzenia umożliwiające wykonanie tych akcji oraz jest w stanie ACTIVE, WORKING lub PAUSED
    \item[Przypadek 1 - Brak wymaganych urządzeń:]Robot zgłasza błąd i przerywa działanie
        \lstinputlisting{listings/test1-1}

    \item[Przypadek 2 - Robot nie implementuje akcji:]Zgłoszenie ostrzeżenia (domyślnie wysyłany jest komunikat zwrotny o niepowodzeniu przypisania zachowania)
        \lstinputlisting{listings/test1-2}

    \item[Przypadek 3 - Nieprawidłowy stan:]Brak działania (domyślnie dostarczono zły komunikat, więc jest on ignorowany)
        \lstinputlisting{listings/test1-3}

    \item[Przypadek 4 - Warunki spełnione:]Następuje przypisanie zachowania, a\,następnie wygenerowanie zdarzenia BEHAVIOUR\_START, które przechwycone przez klasę {\tt Robot} uruchamia je.
        \lstinputlisting{listings/test1-4}
\end{description}

\subsubsection{\underline{Ostrzeżenie o przeszkodzie}}

\begin{description}
    \item[Warunki:]Robot posiada sensor nacisku oraz zachowanie ma zdefiniowaną reakcję na zdarzenie OBSTACLE\_DETECTED
    \item[Przypadek 1 - Brak sensora:]Jeśli robot powinien go posiadać, zgłaszany jest błąd. Jeśli nie - to zdarzenie nie będzie generowane
        \lstinputlisting{listings/test2-1}
    \item[Przypadek 2 - Brak zdefiniowanej reakcji:]Zdarzenie jest wygenerowane, ale nie podjęte są żadne działania
        \lstinputlisting{listings/test2-2}
    \item[Przypadek 3 - Warunki spełnione:]Zgłoszone zdarzenie jest przesłane do\,zachowania i następuje natychmiastowe zatrzymanie robota, a następnie przejście w stan pośredni oraz kontynuowanie zachowania w kolejnym domyślnym stanie.\\
        \lstinputlisting{listings/test2-3}
\end{description}

\subsubsection*{\underline{Ostrzeżenie o bliskim obiekcie}}

\begin{description}
    \item[Warunki:]Robot posiada sensor odległości oraz zachowanie ma zdefiniowaną reakcję na zdarzenie PROXIMITY\_ALERT
    \item[Przypadek 1 - Brak sensora:]Jeśli robot powinien go posiadać, zgłaszany jest błąd. Jeśli nie - to zdarzenie nie będzie generowane
        \lstinputlisting{listings/test3-1}
    \item[Przypadek 2 - Brak zdefiniowanej reakcji:]Zdarzenie jest wygenerowane, ale nie podjęte są żadne działania
        \lstinputlisting{listings/test3-2}
    \item[Przypadek 3 - Warunki spełnione:]Zgłoszone zdarzenie jest przesłane do\,zachowania i następuje natychmiastowe zatrzymanie robota, a następnie przejście w stan pośredni oraz kontynuowanie zachowania w kolejnym domyślnym stanie.\\
        \lstinputlisting{listings/test3-3}
\end{description}

\subsubsection{\underline{Niski poziom baterii}}

Przetestowany dla poziomu większego niż minimalne dopuszczalne napięcie. Po\,jego przekroczeniu, wygenerowane zostało zdarzenie BATTERY\_LOW, wysłano komunikat o niskim naładowaniu baterii oraz nastąpiło zamknięcie aplikacji.

\subsection{Testowanie przejść między stanami}

Komunikaty niezarządzające przetwarzaniem danego stanu są ignorowane.

\subsubsection{\underline{Stan IDLE}}

\begin{description}
    \item[Początek stanu -]Dioda miga na żółto
    \item[Brak odpowiedzi od nadzorcy -]Po ustalonym czasie, przejście w stan PANIC
        \lstinputlisting{listings/test6-idle-1}
    \item[Synchronizacja z nadzorcą -]Nadanie identyfikatora agentowi i przejście w\,stan ACTIVE
        \lstinputlisting{listings/test6-idle-2}
\end{description}

\subsubsection{\underline{Stan ACTIVE}}

\begin{description}
    \item[Początek stanu -]Dioda miga na zielono
    \item[Brak odpowiedzi od nadzorcy -]Po ustalonym czasie, przejście w stan PANIC
        \lstinputlisting{listings/test6-active-1}
    \item[Przesłanie zachowania i komunikatu START -]Przypisanie zachowania, przejście w stan WORKING, rozpoczęcie przetwarzania zachowania.
        \lstinputlisting{listings/test6-active-2}
\end{description}

\subsubsection{\underline{Stan WORKING}}

\begin{description}
    \item[Początek stanu -]Dioda świeci stale na zielono
    \item[Brak odpowiedzi od nadzorcy -]Po ustalonym czasie, przejście w stan PAUSED oraz zatrzymanie zachowania
        \lstinputlisting{listings/test6-working-1}
    \item[Przesłanie komunikatu STOP -]Zatrzymanie zachowania, przejście do\,stanu ACTIVE
        \lstinputlisting{listings/test6-working-2}
    \item[Przesłanie komunikatu PAUSE -]Zatrzymanie zachowania, przejście do\,stanu PAUSED
        \lstinputlisting{listings/test6-working-3}
\end{description}

\subsubsection{\underline{Stan PAUSED}}

\begin{description}
    \item[Początek stanu -]Dioda świeci stale na pomarańczowo
    \item[Brak odpowiedzi od nadzorcy -]Po ustalonym czasie, przejście w stan PANIC
        \lstinputlisting{listings/test6-paused-1}
    \item[Przesłanie komunikatu RESUME -]Wznowienie zachowania, przejście do\,stanu WORKING
        \lstinputlisting{listings/test6-paused-2}
    \item[Przesłanie komunikatu STOP -]Zatrzymanie zachowania, przejście do\,stanu ACTIVE
        \lstinputlisting{listings/test6-paused-3}
\end{description}

\section{Testy funkcjonalne}

\subsection{Testowanie zachowania złożonego i komunikacji}

Na poniższym listingu przedstawiono jeden pełny przebieg zachowania wraz z komunikatami komunikacyjnymi i przejściami między stanami.

\lstinputlisting{listings/test-explore}


\chapter{Podsumowanie}
\label{ch:podsumowanie}


\nocite{*}
\printbibheading[title={Bibliografia}]
\sloppy
\printbibliography[nottype=online,heading=subbibliography,title={Publikacje}]
\printbibliography[type=online,heading=subbibliography,title={Źródła internetowe}]


% \glsaddall
\printglossary[title=Słownik pojęć,toctitle=Słownik pojęć]

\cleardoublepage
\phantomsection
\addcontentsline{toc}{chapter}{\listfigurename}
\listoffigures

\appendix
\chapter{Załączniki}
    % \section{Konfiguracja systemu}
\label{ch:konfiguracja_systemu}

W tym załączniku znajduje się instrukcja kompilacji aplikacji oraz sposobu uruchamiania.

\subsection{Kompilacja}

\textbf{Uwaga:} Kompilacja konfiguracji na architekturę typu ARM wymaga zainstalowanego programu {\tt arm-linux-gnueabi-g++}. Pozostałe dwie konfiguracje wymagają standardowego {\tt g++}.

\subsubsection{W środowisku graficznym}

Wgrane na płytę CD pliki zawierają projekt aplikacji dla środowiska NetBeans (użyta wersja: 8.0.2). Po wczytaniu projektu, na górze programu można wybrać z\,listy rozwijanej bieżącą konfigurację i ją skompilować (klawisz F11).

\subsubsection{W oknie konsoli}

Projekt aplikacji NetBeans tak naprawdę korzysta z zapisanych przez siebie plików Makefile. W katalogu głównym znajduje się plik ogólny Makefile, a w folderze {\tt nbproject} pliki poszczególnych konfiguracji. Proces budowania można rozpocząć poleceniem:
\begin{center}
    {\tt make <nazwa-konfiguracji>}\  lub\
    {\tt make all},
\end{center}
a usunięcie plików tymczasowych wykonuje komenda: {\tt make clean}

\subsection{Synchronizacja plików}

W celu szybkiego wysłania plików do wielu robotów, przygotowany został specjalny katalog {\tt sync} wraz ze skryptem kopiującym. Edytując skrypt można podać następujące wartości parametrów:
\begin{itemize}
    \item Adresy ip wszystkich docelowych urządzeń, na które mają być wgrane pliki wykonywalne
    \item Nazwę użytkownika\textsuperscript{*}
    \item Katalog docelowy\textsuperscript{*}
    \item Hasło\textsuperscript{*}
\end{itemize}
\textsuperscript{*} - w przypadku, jeśli ten parametr jest identyczny na każdym urządzeniu.

Skrypt wykonuje kopię plików binarnych z katalogu {\tt dist} dla konfiguracji ARM i umieszcza ją w katalogu {\tt sync}. Następnie, przy pomocy polecenia {\tt rsync} z użyciem podanego hasła, zawartość katalogu zdalnego będzie dokładnym odwzorowaniem bieżącego katalogu.

\subsection{Uruchomienie aplikacji}

Aplikacja może być uruchomiona na dwa sposoby: z poziomu interfejsu graficznego klocka centralnego lub przez zdalne uruchomienie za pomocą połączenia ssh. Ten drugi wariant pozwala lepiej kontrolować komunikaty wyświetlane na ekranie, ale bardziej obciąża urządzenie. Z tego powodu nie powinno się uruchamiać wersji mobilnej z ustawieniami logowania na standardowe wyjście, ale jeśli jest to wymagane, należy ograniczyć rodzaj wyświetlanych wiadomości do minimum.

Przesłany plik wykonywalny można uruchomić w dwóch trybach: agent i master. Ponadto, aplikacja przyjmuje następujące opcjonalne argumenty wywołania:
\begin{itemize}
    % \item model - rodzaj uruchomionego na robocie modelu (tylko w przypadku wariantu {\tt agent})
    \item log=<tryb> - sprecyzowanie minimalnego poziomu logowania wiadomości za pomocą jednego z następujących trybów: {\tt none error warning info verbose debug}.
\end{itemize}

Zakończenie programu można przedwcześnie wymusić przez przesłanie sygnału SIGINT za pomocą kombinacji klawiszy CTRL+C.\\

Przykładowe uruchomienie:
\begin{center}
    {\tt ./ev3dev agent log=warning}
\end{center}

    \section{Zawartość płyty CD}
\label{ch:zawartosc_plyty}

Na dołączonej do pracy płycie CD znajdują się następujące katalogi:

\begin{itemize}
    \item \textbf{doc} - zawiera wygenerowaną w programie Doxygen dokumentację w formie plików html oraz pdf
    \item \textbf{include} - zawiera pliki nagłówkowe aplikacji pogrupowane w moduły
    \item \textbf{nbproject} - katalog projektu aplikacji w środowisku NetBeans. Zawiera także pliki Makefile do kompilacji
    \item \textbf{src} - zawiera pliki źródłowe aplikacji, w szczególności plik startowy
    \item \textbf{sync} - katalog służący synchronizacji plików. Zawiera skrypt kopiujący odpowiednie pliki wykonywalne i przesyłający je do podanych urządzeń.
\end{itemize}

    \section{Zdalne debugowanie aplikacji}
\label{app:gdb}

Narzędzie {\tt brickstrap}, którego twórcy biblioteki ev3dev używają głównie do przygotowywania obrazu systemu dla LEGO EV3, może posłużyć również do debugowania aplikacji na klocku. Możliwe jest także utworzenie zdalnego systemu plików do synchronizacji danych pomiędzy dwoma systemami Linux za pomocą {\tt nfs}. Konfiguracja środowiska wirtualnego\cite{gdb} składa się z następujących kroków:

\begin{enumerate}
    \item Instalacja narzędzia {\tt brickstrap}
    \item Stworzenie wirtualnego środowiska na bazie obrazu systemu identycznego z tym zainstalowanym na robocie LEGO
    \item Kompilacja programu z poziomu powłoki maszyny witualnej (z flagą {\tt -g} w\,celu dołączenia symboli debugowania)
    \item Instalacja {\tt gdbserver} po stronie klocka LEGO oraz {\tt gdb} po stronie środowiska wirtualnego
    \item Uruchomienie skompilowanego pliku binarnego za pomocą {\tt gdbserver} po stronie klocka.
\end{enumerate}

Po wykonaniu wszystkich kroków, opisanych szczegółowo na portalu ev3dev\cite{gdb}, można uruchomić {\tt gdb} z parametrem identycznym jak nazwa programu uruchomiona na robocie. W ten sposób stworzona została interaktywna sesja debuggera i za pomocą komend wpisywanych bezpośrednio do terminala można np. ustalać punkty przerwań czy podglądać wartości zmiennych lokalnych. Z racji tego, że program {\tt gdb} uruchamiany jest w środowisku wirtualnym, można czasami natrafić na nieobsługiwane wywołania systemowe. Większość napotykanych błędów lub ostrzeżeń nie powinna sprawiać problemów, ale może ograniczać niektóre funkcjonalności debuggera.


\end{document}
